\section*{Bab \ref{ch:g11:nucleus-radioactivity} --- Inti Atom dan Keradioaktifan}

\begin{solution}{exo:g11:nucleus-radioactivity:1}
${}^{16}_{8}\mathrm{O}$: $8$ p, $8$ n; ${}^{27}_{13}\mathrm{Al}$: $13$ p,
$14$ n; ${}^{60}_{27}\mathrm{Co}$: $27$ p, $33$ n;
${}^{235}_{92}\mathrm{U}$: $92$ p, $143$ n; ${}^{3}_{1}\mathrm{H}$: $1$ p,
$2$ n. (Setiap kali $N = A - Z$.)
\end{solution}

\begin{solution}{exo:g11:nucleus-radioactivity:2}
Ketiga karbonnya ($Z = 6$) saling merupakan \omterm{def:g11:nucleus-radioactivity:notation}{isotop}.
${}^{14}_{6}\mathrm{C}$ dan ${}^{14}_{7}\mathrm{N}$ berbagi $A = 14$ tetapi
tidak berbagi $Z$: keduanya unsur yang berbeda (awan elektronnya berbeda),
sehingga bukan \omterm{def:g11:nucleus-radioactivity:notation}{isotop}.
\end{solution}

\begin{solution}{exo:g11:nucleus-radioactivity:3}
$A = 210 - 4 = 206$, $Z = 84 - 2 = 82$:
${}^{210}_{84}\mathrm{Po} \to {}^{206}_{82}\mathrm{Pb} +
{}^{4}_{2}\mathrm{He}$ --- inti anaknya adalah timbal-206.
\end{solution}

\begin{solution}{exo:g11:nucleus-radioactivity:4}
${}^{60}_{27}\mathrm{Co} \to {}^{60}_{28}\mathrm{Ni} +
{}^{\;0}_{-1}\mathrm{e}$ dan
${}^{131}_{53}\mathrm{I} \to {}^{131}_{54}\mathrm{Xe} +
{}^{\;0}_{-1}\mathrm{e}$: $A$-nya tetap, $Z$-nya naik satu.
\end{solution}

\begin{solution}{exo:g11:nucleus-radioactivity:5}
$16$ hari $= 2\,T_{1/2}$: $800/4 = \qty{200}{MBq}$. $32$ hari
$= 4\,T_{1/2}$: $800/16 = \qty{50}{MBq}$.
\end{solution}

\begin{solution}{exo:g11:nucleus-radioactivity:6}
(a) $\alpha$ (\omterm{def:g10:universal-gravitation:weight}{berat}, sehingga nyaris tidak disimpangkan); (b) $\beta^-$;
(c) $\gamma$ (netral, menembus); (d) $\beta^+$ (\omterm{rem:g11:nucleus-radioactivity:betaplus}{positron}: seringan (b),
dengan muatan yang berlawanan).
\end{solution}

\begin{solution}{exo:g11:nucleus-radioactivity:7}
${}^{222}_{86}\mathrm{Rn} \to {}^{218}_{84}\mathrm{Po} +
{}^{4}_{2}\mathrm{He}$; \
${}^{218}_{84}\mathrm{Po} \to {}^{214}_{82}\mathrm{Pb} +
{}^{4}_{2}\mathrm{He}$; \
${}^{214}_{82}\mathrm{Pb} \to {}^{214}_{83}\mathrm{Bi} +
{}^{\;0}_{-1}\mathrm{e}$.
\end{solution}

\begin{solution}{exo:g11:nucleus-radioactivity:8}
Fluorin-18 memiliki $9$ neutron sedangkan fluorin-19 yang stabil memiliki
$10$: miskin neutron, di bawah pitanya, sehingga $\beta^+$:
${}^{18}_{9}\mathrm{F} \to {}^{18}_{8}\mathrm{O} + {}^{0}_{+1}\mathrm{e}$.
\end{solution}

\begin{solution}{exo:g11:nucleus-radioactivity:9}
$24$ jam $= 4\,T_{1/2}$: $500/16 \approx \qty{31}{MBq}$. Karena
$2^9 = 512$, sembilan \omterm{def:g11:nucleus-radioactivity:halflife}{waktu paruh} membawa $500/512 < \qty{1}{MBq}$: yaitu
sesudah $9 \times 6 = \qty{54}{h}$.
\end{solution}

\begin{solution}{exo:g11:nucleus-radioactivity:10}
$\frac18 = \frac1{2^3}$: tiga \omterm{def:g11:nucleus-radioactivity:halflife}{waktu paruh}, yaitu
$3 \times \num{5700} =
\num{17100}$ tahun. Tulang dinosaurus berumur sekitar
$\num{17500}$ \omterm{def:g11:nucleus-radioactivity:halflife}{waktu paruh}: yang tersisa sebagian $1/2^{17500}$ --- tidak
tersisa satu \omterm{def:g11:fundamental-interactions:ladder}{atom} karbon-14 pun (satu gram karbon hanya memuat sekitar
$\num{5e22}$ \omterm{def:g11:fundamental-interactions:ladder}{atom}).
\end{solution}

\begin{solution}{exo:g11:nucleus-radioactivity:11}
${}^{14}_{6}\mathrm{C}$: $8$ neutron berbanding $6$--$7$ pada karbon yang
stabil, berarti di atas pitanya, sehingga $\beta^-$.
${}^{11}_{6}\mathrm{C}$: $5$ neutron, miskin neutron, sehingga $\beta^+$.
${}^{238}_{92}\mathrm{U}$: $Z = 92 > 83$, terlalu berat, sehingga $\alpha$.
\end{solution}

\begin{solution}{exo:g11:nucleus-radioactivity:12}
Hanya $\alpha$ yang mengubah $A$: $238 - 206 = 32 = 4x$, sehingga $x = 8$
\omterm{def:g11:nucleus-radioactivity:abg}{alfa}. Kedelapannya membuang $16$ dari $Z$: $92 - 16 + y = 82$
memberikan $y = 6$ peluruhan beta-minus.
\end{solution}

\begin{solution}{exo:g11:nucleus-radioactivity:13}
$6.4/0.1 = 64 = 2^{6}$: enam \omterm{def:g11:nucleus-radioactivity:halflife}{waktu paruh}, yaitu
$6 \times 5.3 = 31.8 \approx 32$ tahun di dalam penyimpanan bertudung.
\end{solution}

\begin{solution}{exo:g11:nucleus-radioactivity:14}
Tiap hari: $8000 \times 86400 \approx \num{6.9e8}$ peluruhan. Sepanjang $80$
tahun: $\num{6.9e8} \times 365 \times 80 \approx \num{2e13}$. Kehidupan
selalu berjalan di atas latar ini --- satu \omterm{def:g11:nucleus-radioactivity:activity}{becquerel} adalah satu \omterm{def:g11:fundamental-interactions:ladder}{atom} tiap
sekon dari $\sim 10^{27}$ \omterm{def:g11:fundamental-interactions:ladder}{atom} tubuhnya; \omterm{def:g10:orders-of-magnitude:unit}{satuannya} mungil, dan cacah
\omterm{def:g11:nucleus-radioactivity:activity}{becquerel} yang terdengar menakutkan bisa jadi hanya setara beberapa pisang.
\end{solution}

\begin{solution}{exo:g11:nucleus-radioactivity:15}
$\num{25000} \times 86400 \approx \num{2.2e9}$ peluruhan tiap hari. Sepuluh
tahun adalah $10/432 \approx 2\%$ satu \omterm{def:g11:nucleus-radioactivity:halflife}{waktu paruh}: jauh kurang dari
separuhnya yang meluruh, sehingga elektronik dan debunyalah, bukan
sumbernya, yang memaksa penggantian. Di luar, $\alpha$-nya mati di dalam
udara kamarnya dan selongsongnya; debu yang terhirup memarkir pemancarnya
di dalam paru-paru, tempat setiap $\alpha$ menumpahkan \omterm{def:g10:energy-conservation:energy}{energinya} ke jaringan
yang hidup --- \omterm{rem:g11:nucleus-radioactivity:dose}{dosis} yang besar (dalam \omterm{rem:g11:nucleus-radioactivity:dose}{sievert}) dari \omterm{def:g11:nucleus-radioactivity:activity}{aktivitas} yang
sederhana.
\end{solution}

\begin{solution}{pb:g11:nucleus-radioactivity:1}
\textbf{1.} ${}^{12}_{6}\mathrm{C}$: $6$ p, $6$ n;
${}^{14}_{6}\mathrm{C}$: $6$ p, $8$ n. Kimia hanya melihat awan elektronnya,
yang ditetapkan $Z = 6$: perilakunya serupa.

\textbf{2.} Makan dan bernapas terus-menerus memperbarui karbon tubuh yang
hidup pada perbandingan atmosfernya; pada saat mati asupannya berhenti dan
peluruhannya berjalan tanpa lawan.

\textbf{3.} ${}^{14}_{6}\mathrm{C} \to {}^{14}_{7}\mathrm{N} +
{}^{\;0}_{-1}\mathrm{e}$ ($\beta^-$). Partikel $\beta$ ini lembut: kain
pembungkus yang berlapis-lapis menghentikannya, sehingga hanya sedikit yang
lolos dari muminya --- dan itulah sebabnya yang diukur adalah cuplikan
karbon kecil yang diambil dari bendanya sendiri.

\textbf{4.} $6.8/13.6 = \frac12$: satu \omterm{def:g11:nucleus-radioactivity:halflife}{waktu paruh}. Mumi itu berumur
sekitar \num{5700} tahun.

\textbf{5.} $13.2/13.6 \approx 0.97$: kainnya pada dasarnya modern ---
paling banter berumur beberapa abad. Sebuah pemalsuan.

\textbf{6.} $13.6/2^{10} = 13.6/1024 \approx 0.013$ peluruhan tiap menit
--- satu cacahan tiap $75$ menit tiap gram, yang tenggelam di dalam latar
alaminya. Sepuluh \omterm{def:g11:nucleus-radioactivity:halflife}{waktu paruh}, yaitu sekitar \num{57000} tahun, adalah
cakrawala praktisnya.

\textbf{7.} $92$ proton, $146$ neutron. Tolakan Coulomb bekerja antara
semua pasangan protonnya menyeberangi intinya sementara perekat kuatnya
hanya mengikat tetangga: inti yang berat hanya bertahan dengan kelebihan
neutron, yaitu perekat tanpa tolakan.

\textbf{8.} ${}^{238}_{92}\mathrm{U} \to {}^{234}_{90}\mathrm{Th} +
{}^{4}_{2}\mathrm{He}$.

\textbf{9.} ${}^{234}_{90}\mathrm{Th} \to {}^{234}_{91}\mathrm{Pa} +
{}^{\;0}_{-1}\mathrm{e}$, lalu ${}^{234}_{91}\mathrm{Pa} \to
{}^{234}_{92}\mathrm{U} + {}^{\;0}_{-1}\mathrm{e}$: uranium muncul kembali,
sebagai uranium-234.

\textbf{10.} Untuk $A$: $238 - 206 = 32$, sehingga ada $8$ langkah
$\alpha$; untuk $Z$: $92 - 16 + y = 82$, sehingga ada $6$ langkah
$\beta^-$.

\textbf{11.} $\num{4.5e9}$ tahun adalah satu \omterm{def:g11:nucleus-radioactivity:halflife}{waktu paruh}: tersisa
$\frac12$. Lima miliar tahun lagi $\approx$ satu \omterm{def:g11:nucleus-radioactivity:halflife}{waktu paruh} berikutnya:
sekitar $\frac14$.

\textbf{12.} Kalor itu menggerakkan gunung api, lempeng tektonik (benua
yang hanyut, gempa bumi) dan kalor panas bumi: ilmu bumi permukaan sebuah
planet yang dihangatkan dari dalam.

\textbf{13.} ${}^{241}_{95}\mathrm{Am} \to {}^{237}_{93}\mathrm{Np} +
{}^{4}_{2}\mathrm{He}$.

\textbf{14.} $\num{3.3e4}$ peluruhan tiap sekon;
$\num{33000} \times 86400 \approx \num{2.9e9}$ tiap hari.

\textbf{15.} Tiap $\alpha$ mengionkan udara yang dilintasinya; ionnya
membawa arus yang mungil di antara elektrode kamarnya. Butiran asap
menangkap ionnya, arusnya turun, dan alarmnya pun berbunyi.

\textbf{16.} Partikel $\alpha$-nya membelanjakan seluruh \omterm{def:g10:energy-conservation:energy}{energinya} di dalam
kamarnya dan tidak ada yang lolos dari selongsongnya: pengionannya
terbesar, kebocorannya nol. Sumber $\gamma$ dengan \omterm{def:g11:nucleus-radioactivity:activity}{aktivitas} yang sama
nyaris tidak akan mengionkan udara kamarnya dan justru akan menyinari
seluruh ruangannya.

\textbf{17.} Sepuluh tahun adalah $10/432 \approx 2\%$ satu \omterm{def:g11:nucleus-radioactivity:halflife}{waktu paruh}:
aktivitasnya pada dasarnya tidak berubah. Mata rantai yang lemah adalah
debu, serangga dan elektroniknya --- \omterm{def:g11:fundamental-interactions:ladder}{inti atomnya} hidup lebih lama daripada
alatnya.

\textbf{18.} $\num{9e9}$ tahun $= 2$ \omterm{def:g11:nucleus-radioactivity:halflife}{waktu paruh}: bertahan $\frac14$.
Uranium Bumi masih melimpah, sehingga uranium itu ditempa paling banter
beberapa \omterm{def:g11:nucleus-radioactivity:halflife}{waktu paruh} yang lalu --- planetnya memadat dari abu bintang yang
masih segar, miliaran tahun sesudah bintang yang pertama.

\textbf{19.} Mati secara geologi: bagian dalamnya dingin, tidak ada
kegunungapian atau lempeng tektonik yang mendaur ulang udara dan batuannya,
dan intinya tidak bergolak --- sehingga tidak ada tameng \omterm{def:g11:magnetic-fields:magnet}{magnet} terhadap
angin surya.

\textbf{20.} Setiap \omterm{def:g11:fundamental-interactions:ladder}{atom} di dalam dirimu yang lebih berat daripada helium
ditempa di dalam bintang yang sekarat, dan sisa yang tidak stabil dari
tempaan itu masih menghangatkan tanah tempatmu berdiri: debu bintang, yang
dijaga tetap hidup oleh \omterm{def:g11:nucleus-radioactivity:halflife}{waktu paruhnya} sendiri.
\end{solution}
